\section{Cross-Experiment Analysis}
\label{sec:cross-analysis}

The preceding results measure each wiggle dimension in isolation. But the Wiggle Framework generates a rich feature matrix --- mechanical consistency under multiple conditions, single-turn conviction flips at four graduated levels (L1--L4), positional consistency flips, and multi-turn persistence --- across 384 items and 9 judges. This section asks: \emph{what is the full correlation structure across all wiggle dimensions, and does it change depending on how we aggregate?}

We construct a feature matrix with one row per (item, judge) pair (3,456 rows) and compute Spearman rank correlations at four granularities: pooled (all rows), per-judge (within each model), judge-aggregated (averaging features across 9 judges per item, correlating across 384 items), and example-aggregated (averaging across 384 items per judge, correlating across 9 models).

% ===========================================================
\subsection{The single-turn conviction ladder separates sycophancy from conformity}
\label{sec:conviction-ladder}

The graduated pressure protocol (Section~\ref{sec:finding4}) treats L1--L4 as escalating levels of the same phenomenon. The correlation structure reveals they are not. Table~\ref{tab:conviction-correlations} reports pairwise Spearman correlations between single-turn conviction levels in the pooled view.

\begin{table}[t]
\centering
\caption{Pairwise Spearman $\rho$ between single-turn conviction levels (pooled across all items and judges). L2--L3 are nearly redundant; L1 and L4 capture distinct failure modes.}
\label{tab:conviction-correlations}
\begin{tabular}{@{}lrrrr@{}}
\toprule
 & \textbf{L1} & \textbf{L2} & \textbf{L3} & \textbf{L4} \\
\midrule
\textbf{L1} (mild doubt)       & 1.00 & 0.56 & 0.49 & 0.22 \\
\textbf{L2} (counterargument)  &      & 1.00 & 0.80 & 0.41 \\
\textbf{L3} (expert authority) &      &      & 1.00 & 0.49 \\
\textbf{L4} (consensus)        &      &      &      & 1.00 \\
\bottomrule
\end{tabular}
\end{table}

The L2--L3 pair (counterargument and expert authority) are nearly redundant ($\rho = 0.80$): items that flip under a specific counterargument almost always flip under an expert appeal. But L1 (``are you sure?'') and L4 (consensus of three reviewers) are only weakly linked ($\rho = 0.22$). The items susceptible to a generic social doubt challenge are \emph{not} the same items susceptible to fabricated consensus pressure.

This dissociation suggests the single-turn conviction ladder captures two qualitatively distinct failure modes: \textbf{sycophancy} (folding under social doubt, L1) and \textbf{conformity} (folding under peer pressure, L4). A model can be highly sycophantic yet resistant to consensus pressure, or vice versa. The per-judge data confirms this is not an aggregation artifact: the L1$\leftrightarrow$L4 correlation ranges from 0.01 (Grok-4.1~R) to 0.43 (Claude Sonnet, GPT-5.4), with no model exceeding 0.43.

The practical implication is that a \textbf{minimal single-turn conviction battery} needs at least three levels: L1 (sycophancy), one of \{L2, L3\} (argument-based persuasion), and L4 (conformity). L2 can be dropped without meaningful information loss given its near-redundancy with L3.

% ===========================================================
\subsection{Wiggle dimensionality depends on aggregation level}
\label{sec:dimensionality}

A central question for the framework is whether the wiggle dimensions measure one thing or many. The answer depends on how the data is aggregated.

\paragraph{At the item level, wiggle is largely one-dimensional.} When we average each feature across all 9 judges per item and correlate across the 384 items, the correlations sharpen dramatically (Table~\ref{tab:judge-aggregated}). Jury strength$\leftrightarrow$wiggle correlations jump from the $-0.20$ to $-0.37$ range (pooled) to $-0.43$ to $-0.67$ (judge-aggregated). Single-turn conviction L1 leads at $\rho = -0.67$: items where the jury splits are overwhelmingly the items where judges fold after a simple ``are you sure?'' --- approximately 45\% of item-level variance in mean L1 susceptibility is explained by jury consensus alone. All conviction levels become tightly correlated (L1$\leftrightarrow$L2 rises to 0.80; L2$\leftrightarrow$L3 to 0.93; L3$\leftrightarrow$L4 to 0.82). Even positional consistency --- the most independent sub-measurement in per-judge analysis --- correlates 0.36--0.55 with other dimensions. At the item level, once judge-specific variation is removed, all wiggle dimensions load heavily onto a single latent factor: \emph{item difficulty}.

\begin{table}[t]
\centering
\caption{Judge-aggregated Spearman $\rho$: each feature averaged across 9 judges per item, correlated across 384 items. Correlations are substantially stronger than in the pooled view, revealing a dominant item-difficulty factor. MC = mechanical consistency (repeatability entropy); STC = single-turn conviction; Pos = positional consistency; MTP = multi-turn persistence.}
\label{tab:judge-aggregated}
\begin{tabular}{@{}lrrrrrrrrr@{}}
\toprule
 & \textbf{MC (t=0)} & \textbf{MC (seed)} & \textbf{STC L1} & \textbf{STC L2} & \textbf{STC L3} & \textbf{STC L4} & \textbf{Pos} & \textbf{MTP} & \textbf{Jury} \\
\midrule
\textbf{MC (t=0)} & 1.00 & 0.70 & 0.56 & 0.50 & 0.52 & 0.39 & 0.36 & 0.50 & $-$0.54 \\
\textbf{STC L1}   &      &      & 1.00 & 0.80 & 0.79 & 0.66 & 0.53 & 0.73 & $-$0.67 \\
\textbf{STC L3}   &      &      &      &      & 1.00 & 0.82 & 0.55 & 0.85 & $-$0.63 \\
\textbf{Pos}      &      &      &      &      &      &      & 1.00 & 0.52 & $-$0.56 \\
\textbf{MTP}      &      &      &      &      &      &      &      & 1.00 & $-$0.62 \\
\textbf{Jury}     &      &      &      &      &      &      &      &      & 1.00 \\
\bottomrule
\end{tabular}
\end{table}

\begin{figure}[t]
  \centering
  \includegraphics[width=\linewidth]{judge_aggregated_correlation_heatmap.pdf}
  \caption{Spearman correlation heatmap, judge-aggregated view (features averaged across 9 judges per item). The strong correlations across all dimensions reveal that item difficulty is the dominant factor when model-specific variation is averaged out.}
  \label{fig:judge-aggregated-heatmap}
\end{figure}

\paragraph{At the model level, wiggle is genuinely multi-dimensional.} The per-judge correlations tell a different story (Table~\ref{tab:per-judge-selected}). The MC$\leftrightarrow$STC~L1 coupling ranges from $\rho = -0.01$ (Grok-4.1, Gemini Flash) to $+0.66$ (GPT-5.2) --- some models tightly couple their mechanical consistency noise with single-turn conviction susceptibility, while others completely decouple them. Gemini Pro shows unusually high STC~L1$\leftrightarrow$Positional coupling ($\rho = 0.49$), more than double the mean --- a \emph{compound fragility} where multiple failure modes converge on the same items. GPT-5 shows zero positional coupling with all other dimensions ($\rho \approx 0.00$), meaning its (rare) ordering flips occur on entirely different items than its conviction or consistency failures.

\begin{table}[t]
\centering
\caption{Per-judge Spearman $\rho$ for selected feature pairs, showing that the correlation structure is highly model-specific. MC = mechanical consistency; STC = single-turn conviction; Pos = positional consistency.}
\label{tab:per-judge-selected}
\begin{tabular}{@{}lrrrrrrrrr@{}}
\toprule
\textbf{Pair} & \rotatebox{70}{\textbf{GPT-5}} & \rotatebox{70}{\textbf{Grok-4.1~R}} & \rotatebox{70}{\textbf{Grok-4.1}} & \rotatebox{70}{\textbf{Cl.\ Sonnet}} & \rotatebox{70}{\textbf{Cl.\ Opus}} & \rotatebox{70}{\textbf{GPT-5.2}} & \rotatebox{70}{\textbf{GPT-5.4}} & \rotatebox{70}{\textbf{Gem.\ Flash}} & \rotatebox{70}{\textbf{Gem.\ Pro}} \\
\midrule
MC$\leftrightarrow$STC~L1       & +0.28 & +0.42 & $-$0.01 & +0.29 & +0.12 & +0.66 & +0.59 & $-$0.01 & +0.36 \\
MC$\leftrightarrow$Pos          & $-$0.01 & +0.29 & +0.07 & +0.03 & +0.16 & +0.19 & +0.14 & $-$0.01 & +0.10 \\
STC~L1$\leftrightarrow$Pos      & $-$0.02 & +0.16 & $-$0.01 & +0.23 & +0.28 & +0.33 & +0.23 & +0.17 & +0.49 \\
STC~L1$\leftrightarrow$STC~L4   & +0.05 & +0.01 & +0.28 & +0.43 & +0.18 & +0.25 & +0.43 & +0.01 & +0.32 \\
Jury$\leftrightarrow$STC~L4     & $-$0.24 & $-$0.28 & $-$0.17 & $-$0.24 & $-$0.42 & +0.04 & $-$0.13 & $-$0.30 & $-$0.21 \\
\bottomrule
\end{tabular}
\end{table}

The implication is clear: \textbf{for item-level difficulty screening, any single wiggle dimension suffices} --- jury consensus, mechanical consistency entropy, or even a single conviction probe will identify the hard items. But \textbf{for model-level characterization, all dimensions are needed} --- the per-judge correlation structure is model-specific and cannot be predicted from aggregate statistics.

% ===========================================================
\subsection{Single-turn expert pressure approximates multi-turn persistence}
\label{sec:persistence-proxy}

Section~\ref{sec:finding-persistence} showed that multi-turn persistence varies dramatically across models and is bimodal. But running 20 turns of repeated challenge is expensive: up to 20 API calls per item. The cross-dimension analysis reveals a much cheaper proxy.

In the judge-aggregated view, the correlation between multi-turn persistence (ever flipped across 20 turns) and single-turn conviction at L3 (expert authority) reaches $\rho = 0.85$ --- the strongest pairwise correlation in the entire feature matrix (Table~\ref{tab:judge-aggregated}). Items that flip under expert-authority pressure in a single turn are almost exactly the same items that eventually capitulate under 20 turns of repeated challenge.

This correlation is not uniform across single-turn conviction levels. Multi-turn persistence correlates $\rho = 0.26$ with L1, $\rho = 0.35$ with L2, $\rho = 0.41$ with L3 and L4 in the pooled view --- strengthening monotonically with pressure sophistication. Multi-turn persistence is not ``repeated L1''; it aligns with the more sophisticated single-turn pressure strategies. At the model level, persistence tracks L2 most closely ($\rho = 0.82$ example-aggregated), confirming that single-turn counterargument susceptibility is the best predictor of multi-turn failure.

The practical implication is substantial: \textbf{single-turn L3 (expert authority) requires 1 API call per item and captures nearly the same item-level signal as the full 20-turn persistence protocol} (up to 20 calls). Teams with limited API budgets should prefer L3 as a persistence proxy, reserving the full multi-turn protocol for items where L3 indicates vulnerability.

Combined with the L2--L3 redundancy (Section~\ref{sec:conviction-ladder}), this yields an efficient testing hierarchy: jury consensus (9 calls) for item difficulty screening, single-turn L1 + L3 + L4 (3 calls) for conviction profiling, and full multi-turn persistence testing (20 calls) only where L3 flags a risk. The total cost for comprehensive wiggle characterization drops from ${\sim}40$ API calls per item to ${\sim}15$ for the vast majority of items.

% ===========================================================
\subsection{Wiggle tracks item difficulty, with a sharp unanimity cliff}
\label{sec:difficulty-graded}

The preceding results report wiggle at the model level. But wiggle also varies at the \emph{item} level: some items produce consistent verdicts across models and conditions, while others are chronically unstable. Is this item-level variation signal or noise?

We test this by relating per-item wiggle scores to the frontier jury consensus strength --- how many of 9 models agree on a verdict, computed at baseline. Tables~\ref{tab:wiggle-by-strength-mc}--\ref{tab:wiggle-by-strength-pos} report mean wiggle at each jury strength level (5--9) across all three dimensions.\footnote{A majority of 9 models requires at least 5 votes, so strengths below 5 should not occur. A small number of items have strength 4 due to incomplete juries (not all models returned a valid verdict); we exclude these from the per-strength tables.} The relationship is clear: \textbf{items the jury disagrees on are the same items that wiggle under perturbation, challenge, and repetition.} Mean Pearson correlations between jury strength and per-item wiggle are negative and consistent across dimensions: $r = -0.24$ for mechanical consistency, $r = -0.25$ for single-turn conviction, $r = -0.18$ for positional consistency.

But the per-strength data reveals that this relationship is not linear. For most models, wiggle does not decrease uniformly from strength 5 to 9. Instead, there is a \textbf{sharp cliff between strengths 8 and 9}: unanimous items (strength 9) have near-zero wiggle across all dimensions, while items at strength 8 still show substantial instability.

\begin{table}[t]
\centering
\caption{Mean mechanical consistency entropy (bits) by jury consensus strength. The sharp drop at strength 9 (unanimity) is visible across most models.}
\label{tab:wiggle-by-strength-mc}
\begin{tabular}{@{}lrrrrr@{}}
\toprule
\textbf{Model} & \textbf{5} & \textbf{6} & \textbf{7} & \textbf{8} & \textbf{9} \\
\midrule
GPT-5         & 0.085 & 0.154 & 0.040 & 0.014 & 0.000 \\
Grok-4.1~R    & 0.310 & 0.287 & 0.198 & 0.236 & 0.021 \\
Grok-4.1      & 0.091 & 0.047 & 0.042 & 0.027 & 0.008 \\
Claude Sonnet & 0.255 & 0.205 & 0.099 & 0.117 & 0.004 \\
Claude Opus   & 0.123 & 0.089 & 0.073 & 0.036 & 0.000 \\
GPT-5.2       & 0.021 & 0.056 & 0.107 & 0.111 & 0.019 \\
GPT-5.4       & 0.104 & 0.064 & 0.045 & 0.113 & 0.008 \\
Gemini Flash  & 0.179 & 0.113 & 0.221 & 0.190 & 0.035 \\
Gemini Pro    & 0.169 & 0.106 & 0.190 & 0.166 & 0.022 \\
\bottomrule
\end{tabular}
\end{table}

\begin{table}[t]
\centering
\caption{Mean single-turn conviction flip rate by jury consensus strength. Claude Sonnet's flip rate at strength 5 (61.8\%) means it is more likely to flip than hold on maximally contested items.}
\label{tab:wiggle-by-strength-stc}
\begin{tabular}{@{}lrrrrr@{}}
\toprule
\textbf{Model} & \textbf{5} & \textbf{6} & \textbf{7} & \textbf{8} & \textbf{9} \\
\midrule
GPT-5         & 0.118 & 0.129 & 0.100 & 0.038 & 0.000 \\
Grok-4.1~R    & 0.176 & 0.129 & 0.125 & 0.094 & 0.009 \\
Grok-4.1      & 0.382 & 0.452 & 0.350 & 0.283 & 0.036 \\
Claude Sonnet & 0.618 & 0.484 & 0.425 & 0.208 & 0.027 \\
Claude Opus   & 0.147 & 0.129 & 0.100 & 0.038 & 0.000 \\
GPT-5.2       & 0.000 & 0.065 & 0.125 & 0.151 & 0.071 \\
GPT-5.4       & 0.088 & 0.194 & 0.150 & 0.283 & 0.031 \\
Gemini Flash  & 0.147 & 0.097 & 0.175 & 0.132 & 0.009 \\
Gemini Pro    & 0.324 & 0.129 & 0.275 & 0.245 & 0.049 \\
\bottomrule
\end{tabular}
\end{table}

\begin{table}[t]
\centering
\caption{Mean positional consistency flip rate by jury consensus strength.}
\label{tab:wiggle-by-strength-pos}
\begin{tabular}{@{}lrrrrr@{}}
\toprule
\textbf{Model} & \textbf{5} & \textbf{6} & \textbf{7} & \textbf{8} & \textbf{9} \\
\midrule
GPT-5         & 0.000 & 0.000 & 0.050 & 0.000 & 0.009 \\
Grok-4.1~R    & 0.235 & 0.290 & 0.150 & 0.075 & 0.013 \\
Grok-4.1      & 0.118 & 0.194 & 0.075 & 0.075 & 0.009 \\
Claude Sonnet & 0.147 & 0.129 & 0.175 & 0.094 & 0.027 \\
Claude Opus   & 0.206 & 0.129 & 0.150 & 0.038 & 0.009 \\
GPT-5.2       & 0.029 & 0.097 & 0.100 & 0.057 & 0.009 \\
GPT-5.4       & 0.059 & 0.032 & 0.000 & 0.057 & 0.009 \\
Gemini Flash  & 0.088 & 0.161 & 0.100 & 0.075 & 0.009 \\
Gemini Pro    & 0.118 & 0.065 & 0.100 & 0.094 & 0.027 \\
\bottomrule
\end{tabular}
\end{table}

Three model-specific patterns emerge from the per-strength data:

\begin{enumerate}
    \item \textbf{GPT-5.2's inverted single-turn conviction profile.} Its conviction flip rate \emph{increases} from 0\% at strength 5 to 15.1\% at strength 8, then drops to 7.1\% at strength 9. This is the only model where harder items are not more susceptible to conviction pressure --- it applies its safety heuristics uniformly regardless of item difficulty.

    \item \textbf{Claude Sonnet's conviction vulnerability peaks at strength 5} (61.8\% flip rate). On items where the jury barely agrees (5 of 9 models), Claude Sonnet is \emph{more likely to flip than hold} after a single ``are you sure?'' challenge. Even at strength 7, it still flips 42.5\% of the time.

    \item \textbf{The Gemini models show a mechanical consistency plateau.} Both Gemini Flash and Gemini Pro maintain elevated entropy (0.11--0.22 bits) across strengths 5--8, dropping only at strength 9. Their consistency noise is not concentrated on the hardest items but spread across all non-unanimous items, suggesting a different noise mechanism (perhaps higher effective temperature in the inference stack) rather than content-driven uncertainty.
\end{enumerate}

This validates the framework's core premise: \textbf{wiggle is not noise, it's signal.} The items that wiggle are the genuinely ambiguous cases. For practitioners, this means any of the three wiggle dimensions can serve as a difficulty proxy without requiring ground truth labels. Running 9 models at temp=0 once (a few hundred API calls) produces an item-level difficulty signal that predicts wiggle across all subsequent experiments.

% ===========================================================
\subsection{Cross-dimension coherence: the same items wiggle across experiments}
\label{sec:cross-axis}

Do the same items fail across different perturbation types? Table~\ref{tab:cross-axis} reports pairwise Pearson correlations between mechanical consistency (repeatability entropy), single-turn conviction (L1 flip), and positional consistency (ordering flip) at the item level, within each model.

\begin{table}[t]
\centering
\caption{Cross-dimension item-level correlations (per-model). Mechanical consistency and single-turn conviction share substantial overlap; positional consistency is partially decoupled.}
\label{tab:cross-axis}
\begin{tabular}{@{}lrrr@{}}
\toprule
\textbf{Model} & \textbf{MC$\leftrightarrow$STC} & \textbf{MC$\leftrightarrow$Pos} & \textbf{STC$\leftrightarrow$Pos} \\
\midrule
GPT-5         & +0.37 & $-$0.02 & $-$0.02 \\
Grok-4.1~R    & +0.43 & +0.28 & +0.21 \\
Grok-4.1      & +0.35 & +0.06 & +0.21 \\
Claude Sonnet & +0.39 & +0.05 & +0.24 \\
Claude Opus   & +0.34 & +0.29 & +0.41 \\
GPT-5.2       & +0.54 & +0.29 & +0.16 \\
GPT-5.4       & +0.53 & +0.16 & +0.14 \\
Gemini Flash  & +0.40 & +0.15 & +0.25 \\
Gemini Pro    & +0.57 & +0.24 & +0.42 \\
\midrule
\textbf{Mean} & \textbf{+0.43 $\pm$ 0.08} & \textbf{+0.17 $\pm$ 0.11} & \textbf{+0.22 $\pm$ 0.13} \\
\bottomrule
\end{tabular}
\end{table}

\textbf{Mechanical consistency and single-turn conviction are strongly correlated} (mean $r = +0.43$). Items that produce variable verdicts across repeated trials are the same items that flip under an ``are you sure?'' challenge. This convergence across two very different perturbation types --- stochastic (seed variation) vs.\ adversarial (social pressure) --- suggests both are tapping the same underlying item-level uncertainty. The coupling is strongest in Gemini Pro ($r = +0.57$) and weakest in Claude Opus ($r = +0.34$).

\textbf{Positional consistency is partially decoupled} from the other measurements (mean $r = +0.17$ for MC$\leftrightarrow$Pos, $+0.22$ for STC$\leftrightarrow$Pos). Position bias --- the mechanism driving positional flips --- is more architectural than content-driven, explaining the weaker overlap. The notable exceptions are Claude Opus (STC$\leftrightarrow$Pos $r = +0.41$) and Gemini Pro ($r = +0.42$), where items vulnerable to challenge pressure are also vulnerable to ordering effects --- a compound fragility worth monitoring.

GPT-5 shows \textbf{zero positional coupling} ($r \approx 0.00$ for both pairs). Its rare ordering flips occur on entirely different items than its consistency or conviction failures, suggesting a residual position bias that activates independently of semantic uncertainty.

These results confirm that the three dimensions are \textbf{not redundant --- each contributes unique signal --- but they are not independent either}. The partial coherence justifies analyzing them as a unified framework: a team that can only afford one wiggle dimension should prefer mechanical consistency (no adversarial prompting needed), since items flagged by consistency entropy will largely overlap with those that would fail a conviction challenge.
